\documentclass[11pt, a4paper]{report}

% --- PACKAGES ---
\usepackage[utf8]{inputenc}
\usepackage[T1]{fontenc}
\usepackage[french]{babel}
\usepackage{geometry}
\geometry{margin=2.5cm}
\usepackage{xcolor}
\usepackage{hyperref}
\usepackage{enumitem}
\usepackage{listings}
\usepackage{listingsutf8}
\usepackage[most]{tcolorbox}
\usepackage{tabularx}
\usepackage{booktabs}

% --- COULEURS ---
\definecolor{primary}{rgb}{0.1, 0.3, 0.5}
\definecolor{bgcode}{rgb}{0.96, 0.96, 0.96}
\definecolor{keyword}{rgb}{0.1, 0.1, 0.7}
\definecolor{comment}{rgb}{0.1, 0.5, 0.1}

% --- LISTINGS ---
\lstset{
    backgroundcolor=\color{bgcode},
    basicstyle=\ttfamily\small,
    keywordstyle=\color{keyword}\bfseries,
    commentstyle=\color{comment}\itshape,
    frame=single,
    rulecolor=\color{lightgray},
    breaklines=true,
    showstringspaces=false,
    inputencoding=utf8,
    extendedchars=true
}

\newtcolorbox{notebox}[1][]{
    colback=blue!5,
    colframe=blue!50!black,
    fonttitle=\bfseries,
    title=#1,
    arc=3mm
}

\title{\textbf{Cheat Sheet DevSecOps}\\Architecture, Nommage et Commandes (C/C++, Python, Java)}
\author{Référentiel CI/CD \& Qualité de Code}
\date{\today}

\begin{document}

\maketitle
\tableofcontents
\newpage

% =========================================================================
\chapter{Architecture et Structure des Projets}
% =========================================================================
L'organisation des fichiers doit respecter le principe du \textbf{Mirroring} (effet miroir). Le dossier des tests doit reproduire exactement l'arborescence du code source.

\section{Structure Standard C, C++ et Python}

\begin{lstlisting}[language=bash, numbers=none]
mon-projet/
|-- src/                # Code source de production
|   |-- api/
|   |   \-- auth.py     # ou auth.c / auth.cpp
|   \-- utils/
|       \-- math.py
\-- tests/              # Dossier dedie aux tests (DETECTE PAR LA CI)
    |-- api/
    |   \-- test_auth.py
    \-- utils/
        \-- test_math.py
\end{lstlisting}

\section{Structure Standard Java (Maven / Gradle)}
Java impose une arborescence très stricte basée sur les conventions Maven.

\begin{lstlisting}[language=bash, numbers=none]
mon-projet-java/
|-- src/
|   |-- main/
|   |   \-- java/
|   |       \-- com/monentreprise/UserService.java
|   \-- test/           # Dossier dedie aux tests Java
|       \-- java/
|           \-- com/monentreprise/UserServiceTest.java
\-- pom.xml             # Fichier de configuration Maven
\end{lstlisting}

% =========================================================================
\chapter{Règles de Nommage (Production et Tests)}
% =========================================================================

\begin{table}[h]
\centering
\renewcommand{\arraystretch}{1.3}
\begin{tabularx}{\textwidth}{|l|X|X|}
\hline
\textbf{Langage} & \textbf{Code Source (Variables, Classes)} & \textbf{Fichiers de Tests (Imposé par la CI)} \\
\hline
\textbf{C / C++} & 
\texttt{snake\_case} pour variables/fonctions. \texttt{MAJUSCULES} pour macros et enums. & 
\texttt{test\_*.c} ou \texttt{*\_test.cpp} \newline \textit{Ex: test\_calcul.c} \\
\hline
\textbf{Python} & 
\texttt{snake\_case} pour variables/fonctions. \texttt{PascalCase} pour les classes. & 
\texttt{test\_*.py} ou \texttt{*\_test.py} \newline \textit{Ex: test\_auth.py} \\
\hline
\textbf{Java} & 
\texttt{camelCase} pour variables/méthodes. \texttt{PascalCase} pour les classes. & 
\texttt{*Test.java}, \texttt{Test*.java}, \texttt{*TestCase.java} \newline \textit{Ex: UserServiceTest.java} \\
\hline
\end{tabularx}
\end{table}

\newpage

% =========================================================================
\chapter{Dictionnaire des Commandes}
% =========================================================================

% =========================================================================
\section{Dictionnaire des Commandes : C et C++}
% =========================================================================

\subsection{Compilation et Vérification (GCC / Clang)}
\begin{itemize}[leftmargin=1.5cm]
    \item \texttt{gcc} / \texttt{clang} : Compilateurs C.
    \item \texttt{g++} / \texttt{clang++} : Compilateurs C++.
    \item \textbf{Options clés :}
    \begin{itemize}
        \item \texttt{-std=c11} ou \texttt{-std=c++17} : Force la norme du langage.
        \item \texttt{-Wall -Wextra} : Active tous les avertissements (warnings) majeurs.
        \item \texttt{-Werror} : Transforme tous les warnings en erreurs (bloque la compilation).
        \item \texttt{-Wshadow} : Alerte si une variable locale masque une variable globale.
        \item \texttt{-O2} : Niveau d'optimisation standard pour la production.
        \item \texttt{-g} : Inclut les symboles de débogage (utile pour Valgrind ou GDB).
        \item \texttt{-fsanitize=address} : Active ASan (détecte les dépassements de mémoire au runtime).
    \end{itemize}
\end{itemize}

\subsection{Formatage Automatique : \texttt{clang-format}}
Assure un style de code homogène sans débat humain.
\begin{itemize}[leftmargin=1.5cm]
    \item \texttt{-i} : Modifie les fichiers directement (in-place). Sans cette option, il affiche juste le résultat dans le terminal.
    \item \texttt{--style=LLVM} (ou GNU, Google) : Applique les règles de formatage spécifiques d'une norme.
\end{itemize}

\subsection{Analyse Statique (SAST \& Qualité) : \texttt{cppcheck} \& \texttt{flawfinder}}
\begin{itemize}[leftmargin=1.5cm]
    \item \texttt{cppcheck} : Linter C/C++.
    \begin{itemize}
        \item \texttt{--enable=all} : Active toutes les vérifications (style, perfs, portabilité).
        \item \texttt{--inconclusive} : Autorise l'outil à remonter des erreurs probables, même s'il n'est pas sûr à 100\%.
        \item \texttt{--xml} : Exporte le résultat au format XML (nécessaire pour SonarQube).
        \item \texttt{--suppress=missingIncludeSystem} : Ignore les erreurs liées aux entêtes système (ex: \texttt{<stdio.h>}).
    \end{itemize}
    \item \texttt{flawfinder} : Scanneur de vulnérabilités de sécurité (C/C++).
    \begin{itemize}
        \item \texttt{--singleline} : Affiche les vulnérabilités sur une seule ligne (plus facile à lire dans les logs CI).
    \end{itemize}
\end{itemize}

\subsection{Couverture de Code : \texttt{gcovr}}
Extrait les données de couverture générées lors de l'exécution des tests.
\begin{itemize}[leftmargin=1.5cm]
    \item \texttt{--xml-pretty} : Génère un fichier XML lisible (format Cobertura pour SonarQube).
    \item \texttt{--print-summary} : Affiche un résumé rapide dans le terminal de la CI.
    \item \texttt{-o coverage.xml} : Définit le nom du fichier de sortie.
\end{itemize}

\newpage

% =========================================================================
\section{Dictionnaire des Commandes : Python}
% =========================================================================

\subsection{Formatage et Linting : \texttt{ruff} (Remplace Pylint/Black)}
Ruff est un linter et formateur ultra-rapide écrit en Rust, utilisé dans notre pipeline moderne.
\begin{itemize}[leftmargin=1.5cm]
    \item \texttt{python -m ruff format .} : Formate automatiquement le code.
    \item \texttt{python -m ruff format --check .} : Vérifie si le code est bien formaté (Renvoie une erreur "exit 1" si ce n'est pas le cas, utile en CI).
    \item \texttt{python -m ruff check .} : Lance l'analyse statique (détecte les variables inutilisées, les imports morts, etc.).
\end{itemize}

\subsection{Typage Statique : \texttt{mypy}}
\begin{itemize}[leftmargin=1.5cm]
    \item \texttt{python -m mypy .} : Vérifie les Type Hints (\texttt{x: int}, \texttt{-> str}).
    \item \texttt{--strict} : Option facultative mais recommandée. Force le typage de \textbf{toutes} les fonctions sans exception.
\end{itemize}

\subsection{Sécurité SAST : \texttt{bandit}}
Analyse le code source Python pour trouver des failles de sécurité.
\begin{itemize}[leftmargin=1.5cm]
    \item \texttt{-r .} : Scanne le dossier courant récursivement.
    \item \texttt{-f xml -o bandit-report.xml} : Spécifie le format de sortie (\texttt{xml}) et le fichier cible pour SonarQube.
\end{itemize}

\subsection{Tests et Couverture : \texttt{pytest} \& \texttt{coverage}}
\begin{itemize}[leftmargin=1.5cm]
    \item \texttt{python -m coverage run -m pytest} : Lance les tests \texttt{pytest} tout en enregistrant quelles lignes de code ont été exécutées.
    \item \texttt{-q} (pour pytest) : Mode "quiet" (silencieux), réduit la verbosité dans la CI.
    \item \texttt{--junitxml=pytest-report.xml} : Exporte le résultat (succès/échecs des tests) en format standardisé XML.
    \item \texttt{python -m coverage xml} : Transforme les données de couverture enregistrées en un fichier \texttt{coverage.xml} lisible par SonarQube.
\end{itemize}

\newpage

% =========================================================================
\section{Dictionnaire des Commandes : Java}
% =========================================================================

\subsection{Compilation et Tests : Maven (\texttt{mvn})}
Maven est l'outil de gestion de build standard.
\begin{itemize}[leftmargin=1.5cm]
    \item \texttt{mvn clean} : Supprime le dossier \texttt{target/} (les anciens builds) pour repartir de zéro.
    \item \texttt{mvn test} : Compile le projet et lance les tests unitaires via le plugin Surefire (qui cherche les fichiers \texttt{*Test.java}).
    \item \texttt{-fn} (Fail Never) : Option cruciale en CI non-bloquante. Même si un test échoue, Maven ne retourne pas d'erreur fatale, permettant au pipeline de continuer et d'envoyer le rapport à SonarQube.
\end{itemize}

\subsection{Couverture de Code : JaCoCo}
\begin{itemize}[leftmargin=1.5cm]
    \item \texttt{mvn jacoco:report} : Génère le rapport de couverture XML dans le dossier \texttt{target/site/jacoco/jacoco.xml}. Doit être lancé \textbf{après} l'exécution des tests.
\end{itemize}

\subsection{Analyse SonarQube et SAST}
En Java, SonarQube fait souvent l'analyse statique lui-même, mais il peut être complété.
\begin{itemize}[leftmargin=1.5cm]
    \item \texttt{mvn sonar:sonar} : Envoie le code, les rapports de test (Surefire) et les rapports de couverture (JaCoCo) vers le serveur SonarQube.
    \item \texttt{-Dsonar.projectKey=...} : Définit l'identifiant unique du projet sur le serveur Sonar.
    \item \texttt{-Dsonar.host.url=...} : URL du serveur SonarQube (ex: http://localhost:9000).
    \item \texttt{-Dsonar.token=...} : Jeton d'authentification pour se connecter au serveur.
\end{itemize}

\newpage

% =========================================================================
\section{Résumé du Pipeline CI/CD DevSecOps}
% =========================================================================
La logique appliquée à tous ces outils dans notre GitHub Actions repose sur deux règles de base :
\begin{enumerate}
    \item \textbf{Vérification bloquante du Format et du Style} : Si le code est mal nommé (ex: \texttt{test\_*.py}), mal indenté (\texttt{clang-format}, \texttt{Ruff}), la CI s'arrête en erreur (\texttt{exit 1}). Le développeur doit corriger localement et pousser de nouveau.
    \item \textbf{Vérification non-bloquante des Tests (Best Effort)} : Si aucun test n'est trouvé, la CI crée un rapport vide factice (\texttt{touch coverage.xml}) pour ne pas faire planter SonarQube. L'objectif est de mesurer la qualité continue, pas de bloquer le projet s'il n'y a pas encore de tests.
\end{enumerate}

\end{document}